\chapter{Curvature}

\begin{center}\itshape
A flat sheet rolls into a cylinder without changing the lengths drawn on it. It cannot wrap smoothly around a sphere without stretching. Curvature distinguishes the two situations using measurements within the sheet. Added over a closed surface, that local measurement becomes a topological invariant.
\end{center}

In the summers of 1821--25, Gauss worked on the survey of Hanover. Triangles connected his observing stations; his heliotrope reflected sunlight to make distant stations visible. The geometry had to account for the curved Earth. In 1827 he presented the \emph{Disquisitiones generales circa superficies curvas}, with the result he called the \emph{Theorema Egregium}, the remarkable theorem. The curvature calculated from a surface's normal direction can also be calculated from lengths measured on the surface itself.

Gauss measured the large triangle joining the Brocken, Hohehagen and Inselsberg, with sides approximately $69$, $85$ and $107$ kilometres. Historians disagree about the familiar claim that he intended it as a test of Euclidean geometry in physical space: Miller's 1972 criticism and Breitenberger's later account are both listed below. The triangle's spherical angle excess, about fifteen seconds of arc, is explained by the Earth's curvature. Exercise~\ref{ex:girard} computes that scale. The measured triangle and the disputed interpretation of its purpose should be kept separate.

The main calculation uses a moving pair of perpendicular unit vectors. Its connection form records how those directions turn, and its derivative defines curvature. For a surface in space we shall prove that this equals the determinant of the Gauss map. Then Chern's surface argument relates the curvature integral to the indices of zeros of a vector field. A generic height function supplies a field with finitely many zeros, making the resulting Euler characteristic well defined. Finally, on the flat torus, the Laplacian will turn a geometric spectrum into the arithmetic of sums of two squares.

\section{Riemannian metrics}

\begin{definition}
A \emph{Riemannian metric} is an inner product $\langle\ ,\ \rangle_p$ on each tangent space $T_pM$, varying smoothly: in every chart the coefficients $g_{ij}=\langle\partial_{x^i},\partial_{x^j}\rangle$ are smooth. For a piecewise smooth curve $\gamma:[a,b]\to M$, define its length by $\int_a^b|\gamma'(t)|dt$. A diffeomorphism $F$ between Riemannian manifolds is an \emph{isometry} when
$\langle dF_pv,dF_pw\rangle_{F(p)}=\langle v,w\rangle_p$ for all $p$ and $v,w\in T_pM$.
\end{definition}

A submanifold of $\R^N$ inherits the Euclidean inner product on each tangent space. In a chart its squared length element is $\sum_{i,j}g_{ij}dx^i dx^j$. On the unit sphere the usual parametrization by colatitude $\phi$ and longitude $\vartheta$ gives $d\phi^2+\sin^2\phi\,d\vartheta^2$ for $0<\phi<\pi$. The degeneracy of this coordinate expression at the poles is a chart issue, not a degenerate metric.

Every smooth manifold in our convention has a metric. On chart domains take the coordinate Euclidean inner products and weight them by a subordinate partition of unity from Chapter~1. Their locally finite sum is smooth. For each nonzero vector, every defined inner-product term is positive and at least one has a positive weight, so the sum is positive definite. In dimension zero there are no nonzero tangent vectors, and the unique inner products on the zero spaces meet the same definition.

\begin{example}[the hyperbolic half-plane]\label{ex:halfplane}
On $\mathbb H=\{(x,y):y>0\}$, use $(dx^2+dy^2)/y^2$. A vertical segment from height $\varepsilon$ to height $1$, $0<\varepsilon\le1$, has length $\log(1/\varepsilon)$. It takes unbounded length to reach height zero. Real determinant-one fractional linear maps are isometries: for $g(z)=(az+b)/(cz+d)$,
\[
\operatorname{Im}g(z)=\frac{\operatorname{Im}z}{|cz+d|^2},\qquad
|g'(z)|=\frac1{|cz+d|^2},
\]
so the two scale factors cancel in $|dz|/y$. The inverse has the same form. We will compute its curvature as $-1$ here; Chapter~9 develops its distance and isometry group.
\end{example}

\section{Frames and the connection form}

For the surface calculations, let $M$ be oriented and carry a Riemannian metric. Unless a boundary is explicitly mentioned, surfaces in this chapter have no boundary. An \emph{oriented orthonormal frame} $(e_1,e_2)$ on an open set is a smooth pair of positively oriented perpendicular unit tangent fields. Apply Gram--Schmidt to the coordinate vectors of a positively oriented chart to obtain one locally. Its \emph{dual coframe} consists of $1$-forms $\theta^i$ with $\theta^i(e_j)=\delta_{ij}$. The area form is $dA=\theta^1\wedge\theta^2$. A new oriented orthonormal frame differs by a rotation of determinant $1$, so Proposition~\ref{prop:wedgedet} makes this area form independent of the frame.

\begin{proposition}[Cartan's first structure equations]\label{prop:structure}
For every oriented orthonormal coframe, a unique $1$-form $\omega$, called the \emph{connection form}, satisfies
\[
d\theta^1=\omega\wedge\theta^2,\qquad d\theta^2=-\omega\wedge\theta^1.
\]
If the frame changes by $\tilde e_1=\cos\psi\,e_1+\sin\psi\,e_2$, $\tilde e_2=-\sin\psi\,e_1+\cos\psi\,e_2$, with a smooth local angle $\psi$, then $\tilde\omega=\omega+d\psi$.
\end{proposition}

\begin{proof}
Every $2$-form on a surface is a scalar multiple of the area form. Write $d\theta^1=a\theta^1\wedge\theta^2$ and $d\theta^2=b\theta^1\wedge\theta^2$. Then $\omega=a\theta^1+b\theta^2$ satisfies both equations. Their difference for two possible forms would wedge to zero with both coframe members, forcing both of its coefficients to vanish. This proves uniqueness.

The dual coframe changes by $\tilde\theta^1=\cos\psi\,\theta^1+\sin\psi\,\theta^2$ and $\tilde\theta^2=-\sin\psi\,\theta^1+\cos\psi\,\theta^2$. Differentiate the first expression:
\[
d\tilde\theta^1=d\psi\wedge\tilde\theta^2
+\cos\psi\,\omega\wedge\theta^2-\sin\psi\,\omega\wedge\theta^1
=(\omega+d\psi)\wedge\tilde\theta^2.
\]
Differentiating the second gives $d\tilde\theta^2=-(\omega+d\psi)\wedge\tilde\theta^1$. Uniqueness identifies the new form. A smooth local angle exists because the rotation's first column lies in $S^1$ and locally has a smooth argument. A single global angle on the whole overlap is not required.
\end{proof}

\begin{definition}
The \emph{Gauss curvature} $K$ is defined by $d\omega=-K\,dA$.
\end{definition}

Although $\omega$ changes with the frame, its derivative does not: $d(\omega+d\psi)=d\omega$. Thus $K$ is a smooth function determined by the metric. Reversing orientation negates both $\omega$ and $dA$, so it leaves $K$ unchanged as a scalar. The orientation helps the computation but is not part of the answer.

\begin{proposition}[orthogonal coordinates]\label{prop:orthocurv}
In a positively oriented chart with metric $E\,dx^2+G\,dy^2$, put $a=\sqrt E>0$ and $b=\sqrt G>0$. For the coframe $a\,dx,b\,dy$,
\[
\omega=-\frac{a_y}{b}dx+\frac{b_x}{a}dy,\qquad
K=-\frac1{ab}\left[\left(\frac{b_x}{a}\right)_x+\left(\frac{a_y}{b}\right)_y\right].
\]
The scalar curvature formula is unchanged if the coordinate orientation is reversed.
\end{proposition}

\begin{proof}
The two coframe derivatives are $-a_y dx\wedge dy$ and $b_x dx\wedge dy$. For $\omega=p\,dx+q\,dy$, the structure equations require $pb=-a_y$ and $qa=b_x$. These give the claimed connection form. Its derivative is $[(b_x/a)_x+(a_y/b)_y]dx\wedge dy$; compare with $-K ab\,dx\wedge dy$. Reversing orientation changes the signs of the connection and area forms together, preserving $K$.
\end{proof}

\begin{example}\label{ex:curvexamples}
On the sphere, take $a=1$, $b=\sin\phi$ in the colatitude chart. The formula gives $K=-(\cos\phi)'/\sin\phi=1$, and smoothness extends this value to the poles. For the half-plane, $a=b=1/y$, $(a_y/b)_y=(-1/y)_y=1/y^2$, and $K=-1$. For Euclidean polar coordinates, $a=1,b=r$ and $K=-(1)'/r=0$ for $r>0$; the origin is another coordinate singularity, with the same zero curvature in Cartesian coordinates.
\end{example}

\subsection*{Covariant derivative and geodesic curvature}

The connection form differentiates a tangent vector field while accounting for the turning of its frame. For $Y=u_1e_1+u_2e_2$, define
\[
\nabla_vY=(v(u_1)-u_2\omega(v))e_1+(v(u_2)+u_1\omega(v))e_2.
\]
In particular $\nabla e_1=\omega e_2$ and $\nabla e_2=-\omega e_1$. Under a frame rotation, differentiating its coefficients contributes $d\psi$, exactly the extra term in $\tilde\omega$. For example, direct differentiation gives $\nabla\tilde e_1=(\omega+d\psi)\tilde e_2$ and $\nabla\tilde e_2=-(\omega+d\psi)\tilde e_1$; the product rule then proves that the definition is frame independent for every $Y$. This is the \emph{Levi-Civita covariant derivative} on the surface.

Along a curve, replace derivatives $v(u_i)$ by $u_i'(t)$ to define $DY/dt$. The terms involving $\omega$ cancel in the derivative of an inner product, giving
\[
\frac d{dt}\langle Y,Z\rangle
=\left\langle\frac{DY}{dt},Z\right\rangle+
\left\langle Y,\frac{DZ}{dt}\right\rangle.
\]
Thus a vector can turn covariantly without changing its length.

\begin{definition}
For a unit-speed curve with tangent $T$, let $JT$ be the positive right-angle rotation of $T$. Its \emph{geodesic curvature} is the scalar $\kappa_g$ satisfying $DT/ds=\kappa_g JT$; the derivative is perpendicular to $T$ by the metric identity above. A parametrized curve is a \emph{geodesic} when $D\gamma'/dt=0$.
\end{definition}

\begin{proposition}\label{prop:kappag}
If the unit tangent is $T=\cos\varphi\,e_1+\sin\varphi\,e_2$, using a continuous angle along the curve, then $\kappa_g=\varphi'+\omega(T)$.
\end{proposition}

\begin{proof}
Differentiate the two scalar coefficients and the frame. The $e_1$ coefficient of $DT/ds$ is $-(\varphi'+\omega(T))\sin\varphi$, and the $e_2$ coefficient is $(\varphi'+\omega(T))\cos\varphi$. Together they are $(\varphi'+\omega(T))JT$.
\end{proof}

In coordinates the geodesic equation is a smooth system of second-order differential equations. Convert it to a first-order system for position and velocity and apply Theorems~\ref{thm:picard} and~\ref{thm:flowsmooth}. For every initial pair $(p,v)$ there is a unique local geodesic $\gamma_v$, depending smoothly on the data. Its speed is constant because $d|\gamma'|^2/dt=2\langle D\gamma'/dt,\gamma'\rangle=0$. Uniqueness also gives $\gamma_{cv}(t)=\gamma_v(ct)$ wherever both sides are defined.

For small $v$, define the \emph{exponential map} $\exp_p(v)=\gamma_v(1)$. Smooth dependence on a neighborhood of the constant solution at $v=0$ ensures this common time interval for sufficiently small $v$. The scaling identity gives $\exp_p(tv)=\gamma_v(t)$ and hence $d(\exp_p)_0(v)=v$. The inverse function theorem makes $\exp_p$ a diffeomorphism on a small tangent disk. Exercise~\ref{ex:geodlength} explains why its radial geodesics are shortest locally, including curves that try to leave this neighborhood.

The orthogonality needed for local minimization has a short proof here. The structure equations imply that the connection is torsion free:
$\nabla_UV-\nabla_VU=[U,V]$. Indeed, its first coframe coefficient differs from that of $[U,V]$ by $d\theta^1(U,V)-(\omega\wedge\theta^2)(U,V)=0$; the second uses the other structure equation. For a smooth variation $F(r,\vartheta)=\exp_p(r\xi_\vartheta)$ by unit radial geodesics, put $T=F_r$ and $W=F_\vartheta$. Their mixed covariant derivatives agree by this identity, so
\[
\partial_r\langle T,W\rangle
=\langle D_rT,W\rangle+\langle T,D_rW\rangle
=\langle T,D_\vartheta T\rangle
=\tfrac12\partial_\vartheta|T|^2=0.
\]
At $r=0$, $W=0$ and the inner product is zero. Hence the radial and angular directions are perpendicular, while $|T|=1$. This is the \emph{Gauss lemma}: away from the centre, the polar metric is $dr^2+G(r,\vartheta)d\vartheta^2$, with $G>0$ in the normal chart. The variation is smooth even at $r=0$ as a map of $(r,\vartheta)$; only its polar coordinates on the surface fail to be a chart there.

\section{Surfaces in space and the Theorema Egregium}

For an oriented surface $S\subset\R^3$ with its induced metric, choose the normal $N$ so that $(e_1,e_2,N)$ is a positive ambient basis. This defines its \emph{Gauss map} $N:S\to S^2$. Since $|N|=1$, $dN_p$ maps $T_pS$ into $N(p)^\perp$, the same vector plane as $T_pS$. The linear map $-dN_p$ is the \emph{shape operator}.

\begin{theorem}[Gauss, Theorema Egregium]\label{thm:egregium}
For an oriented surface in $\R^3$, $K=\det(dN)$. This determinant, defined using the surface's position in space, depends only on its induced metric and is preserved by isometries.
\end{theorem}

\begin{proof}
On a frame domain put $e_3=N$ and $\omega_{ij}=\langle de_i,e_j\rangle$. Differentiating $\langle e_i,e_j\rangle=\delta_{ij}$ gives $\omega_{ij}=-\omega_{ji}$ and $de_i=\sum_j\omega_{ij}e_j$. The position map $x:S\to\R^3$ satisfies $dx=\theta^1e_1+\theta^2e_2$ because both sides act identically on tangent vectors.

Differentiate this identity componentwise. The product rule is $d(\theta e)=d\theta\,e-\theta\wedge de$ for a $1$-form $\theta$. From $d^2x=0$, the first two vector coefficients give
\[
d\theta^1=\omega_{12}\wedge\theta^2,\qquad
d\theta^2=-\omega_{12}\wedge\theta^1.
\]
Uniqueness of the connection form therefore gives $\omega=\omega_{12}$. In particular, the covariant derivative above is the tangent component of ordinary differentiation in $\R^3$.

Next differentiate $de_1=\omega_{12}e_2+\omega_{13}e_3$ and take its $e_2$ coefficient. The result is $0=d\omega_{12}-\omega_{13}\wedge\omega_{32}$, so
\[
d\omega=-\omega_{31}\wedge\omega_{32}.
\]
But $dN=\omega_{31}e_1+\omega_{32}e_2$. Evaluating its two coefficient forms on $e_1,e_2$ gives its matrix, and Proposition~\ref{prop:wedgedet} gives
$\omega_{31}\wedge\omega_{32}=\det(dN)\,\theta^1\wedge\theta^2$. Compare with $d\omega=-K\,dA$ to obtain the result.
\end{proof}

The normal coefficient in $d^2x=0$ is $\theta^1\wedge\omega_{13}+\theta^2\wedge\omega_{23}=0$. If the matrix of $dN$ has off-diagonal entries $b,c$, this identity reads $(c-b)\theta^1\wedge\theta^2=0$, so it is symmetric. The spectral theorem gives real eigenvalues in perpendicular principal directions. With the convention that principal curvatures are eigenvalues of $-dN$, their product is $K$; changing the normal changes both signs. Normal sections have these signed curvatures, with their extreme values in the principal directions.

For a radius-$R$ sphere with outward normal, $dN=I/R$, so $K=1/R^2$ (the shape-operator eigenvalues are $-1/R$ in this convention). A cylinder has one zero principal curvature and product zero. Cutting a generator lets us unroll it into a planar strip isometrically; without that cut, the whole cylinder is not bijective to a strip. Thus the individual principal curvatures describe bending in space, while their product can be measured within the surface.

\begin{beauty}
The normal direction appears in the proof, but disappears from the answer. The two coframe derivatives determine $\omega$ using the metric. Then $d^2e_1=0$ forces its derivative to agree with the determinant of the normal's derivative. The cancellation in Gauss's theorem is therefore a compatibility condition: a moving orthonormal frame and its normal cannot change independently.
\end{beauty}

\section{The Gauss--Bonnet theorem}

\subsection*{Indices of vector fields}

Near an isolated zero $p$ of a smooth tangent field $X$, choose an oriented chart and a small counterclockwise circle avoiding other zeros. The nonzero coordinate vectors of $X$ along that circle trace a closed curve in $\R^2\setminus\{0\}$. Its winding number is the \emph{index} $\operatorname{ind}_p(X)$. A source turns once positively, a saddle once negatively, and a quadratic field can turn twice.

\begin{lemma}\label{lem:index}
The index is independent of the small circle's radius and the oriented chart. In any oriented orthonormal frame $(f_1,f_2)$ extending over a disk about $p$, write $X/|X|=\cos\psi f_1+\sin\psi f_2$ locally. Then
\[
\operatorname{ind}_p(X)=\frac1{2\pi}\int_c d\psi.
\]
The local angle derivatives fit into a single $1$-form along the loop even when the angle itself does not return to its initial real value.
\end{lemma}

\begin{proof}
Shrinking within an annulus free of zeros gives a homotopy of nonzero vector loops, preserving winding by Volume~7, Proposition~6.6. Under an oriented coordinate change, the component vectors are multiplied by a smooth matrix $A(q)$ with positive determinant. On a small disk centred at $p=0$, replacing $A(q)$ by $A(tq)$, $t$ from $1$ to $0$, deforms this multiplication to the constant matrix $A(0)$ without creating a zero. Write $A(0)=QR$ by Gram--Schmidt, with $Q$ a rotation and $R$ upper triangular with positive diagonal. Rotate $Q$ to the identity and deform $R$ linearly to the identity through such matrices. Thus the constant matrix does not change winding either.

The small loops used by two charts are themselves homotopic with positive orientation in a common punctured neighborhood. For example, linearize the coordinate transition on a sufficiently small circle; its positive-determinant linear image is an ellipse winding once positively, which can be radially deformed to a circle. This justifies comparing the component calculations on the same underlying loop. Frame components are another positive-determinant basis change, so the same argument applies. Finally Corollary~\ref{cor:windingintegral} gives the angle-increment integral, including loops on which only local angle branches exist.
\end{proof}

\begin{figure}[ht]
\centering
\fig{ch04-index}
\caption{Three isolated zeros. Going once around the coral circle anticlockwise, the arrow turns once anticlockwise for the source $(x,y)$, once
clockwise for the saddle $(x,-y)$, and twice anticlockwise for the field $z^2=(x^2-y^2,2xy)$. Their indices are $+1$, $-1$ and $+2$.}
\label{fig:index}
\end{figure}

\begin{theorem}[Gauss--Bonnet--Chern]\label{thm:gaussbonnet}
Let $M$ be a compact oriented Riemannian surface without boundary, and let $X$ be a smooth tangent vector field with finitely many zeros $p_1,\ldots,p_m$. Then
\[
\int_MK\,dA=2\pi\sum_{i=1}^m\operatorname{ind}_{p_i}(X).
\]
\end{theorem}

\begin{proof}
On $M^*=M\setminus\{p_i\}$ the normalized field $e_1=X/|X|$ and $e_2=Je_1$ give a global oriented orthonormal frame. Its connection form satisfies $d\omega=-K\,dA$. Remove disjoint small chart disks $D_i(\varepsilon)$ about the zeros, leaving the compact surface $M_\varepsilon$ with boundary. Its boundary circles are clockwise relative to the disks. Stokes therefore gives
\[
\int_{M_\varepsilon}K\,dA
=-\int_{M_\varepsilon}d\omega
=\sum_i\int_{\partial D_i(\varepsilon)}\omega,
\]
where the circles on the right are counterclockwise.

Choose a smooth local frame $f$ on the full disk, including its centre. Relative to it $e_1$ has a local angle $\psi$, and $\omega=\omega_f+d\psi$. The index lemma and Stokes on the disk give
\[
\int_{\partial D_i(\varepsilon)}\omega
=-\int_{D_i(\varepsilon)}K\,dA+2\pi\operatorname{ind}_{p_i}(X).
\]
The integrals over shrinking disks tend to zero because $K\,dA$ is smooth. Let $\varepsilon\to0$ to obtain the theorem. If there are no zeros, the global frame exists on all of $M$ and Stokes directly gives zero; the empty sum is therefore included. On the empty surface both sides are zero as well.
\end{proof}

The integral does not mention the chosen field, and the index sum does not mention the metric. Thus all fields with finitely many zeros have the same sum, and all metrics have the same curvature integral. A diffeomorphism preserves the oriented index calculation (and the total also survives orientation reversal, which reverses both the domain-loop and vector-coordinate orientations).

To define this common sum for every compact oriented surface, we first construct an admissible field. A smooth function is called \emph{Morse} when each critical point has invertible coordinate Hessian. At a critical point the Hessian changes by $H\mapsto A^{\mathsf T}HA$, so invertibility is independent of coordinates. For a metric, its gradient is the field determined by $\langle\operatorname{grad}f,v\rangle=df(v)$. In coordinates it is $G^{-1}\nabla f$, and has the same zeros as the ordinary partial-derivative vector. Here is an existence proof before we rely on such a function.

Embed $M$ in $\R^N$ by Theorem~\ref{thm:whitney}. Consider the linear height functions $f_a(x)=a\cdot x$. In a graph chart $x=(x',g(x'))$, $x'\in\Omega\subset\R^2$, write $a=(a',a'')$. Such charts exist by Theorem~\ref{thm:fourdesc}(c). Put $F_{a''}(x')=a''\cdot g(x')$; then
\[
\nabla f_a=a'+\nabla F_{a''},\qquad
\operatorname{Hess}f_a=\operatorname{Hess}F_{a''}.
\]
For fixed $a''$, a degenerate critical point requires $-a'$ to be a critical value of the map $\nabla F_{a''}:\Omega\to\R^2$. Equal-dimensional Sard, Theorem~\ref{thm:sard}, makes these $a'$ a null set.

The full exceptional set of $a$ is measurable: it is the image of the Hessian-singular set in $\Omega\times\R^{N-2}$ under the continuous map $(x',a'')\mapsto(-\nabla F_{a''}(x'),a'')$. Exhaust the open domain by compact sets; intersecting each with that relatively closed singular set gives compact sets, so its image is a countable union of compact sets. Fubini now makes the exceptional parameter set null. A countable graph atlas gives a null union of exceptions. Hence almost every $a$ gives a Morse function on all of $M$.

Its critical set is closed. In coordinates at a critical point, the vector of ordinary partial derivatives has invertible Hessian derivative, so the inverse function theorem makes its zero isolated. The Riemannian gradient is obtained by multiplying this vector by the invertible metric matrix inverse, so it has exactly the same zeros. This argument does not require orientation. If it were infinite, compactness would give an accumulation point in that same critical set, contradicting isolation. Thus there are finitely many critical points, exactly the zeros of its gradient for any chosen metric. For $M=\varnothing$, the unique function is Morse vacuously and its field has no zeros.

\begin{definition}
For a compact oriented surface without boundary, its \emph{Euler characteristic} is $\chi(M)=\sum_p\operatorname{ind}_p(X)$ for any smooth tangent field with finitely many zeros.
\end{definition}

The construction shows that this definition applies to every compact oriented surface without boundary. The index-sum theorem shows that the answer is independent of the chosen field and metric.

\begin{corollary}[Gauss--Bonnet]\label{cor:gaussbonnet}
Every Riemannian metric on a compact oriented surface $M$ without boundary satisfies $\int_MK\,dA=2\pi\chi(M)$.
\end{corollary}

\begin{proof}
The generic-height construction just given supplies a smooth field with finitely many zeros. Apply the preceding index-sum theorem to that field.
\end{proof}

\subsection*{Morse functions exist}

A critical point of a smooth real function $f$ has $df=0$. It is \emph{nondegenerate} if its Hessian in local coordinates is invertible. At a critical point a change of coordinates transforms the Hessian by $H\mapsto A^{\mathsf T}HA$; the terms involving first derivatives of $f$ vanish. Thus the condition is coordinate independent. A \emph{Morse function} has only nondegenerate critical points. Its gradient is defined by $\langle\operatorname{grad}f,v\rangle=df(v)$.

\begin{lemma}\label{lem:morseindex}
At a nondegenerate critical point $p$ of a smooth function on a Riemannian surface, $\operatorname{grad}f$ has an isolated zero with index $\operatorname{sgn}\det H$. This index is $+1$ at a local maximum or minimum and $-1$ at a saddle.
\end{lemma}

\begin{proof}
In coordinates centred at $p$, let $G=(g_{ij})$. The gradient's component vector is $G^{-1}\nabla f$, where $\nabla f$ lists the ordinary partial derivatives. The map $\nabla f$ has invertible derivative $H$ at $0$, so the inverse function theorem makes its zero isolated. Multiplying it by the positive definite $G^{-1}$ does not change its zeros.

On a sufficiently small circle, $\nabla f(x)=Hx+o(|x|)$. Since $H$ is invertible there is $c>0$ with $|Hx|\ge c|x|$; the remainder is smaller than this bound, so linear interpolation never passes through zero. This is the winding-number comparison from Volume~7, Proposition~6.6(a). Deform $G^{-1}$ to the identity through positive definite matrices. The index is consequently the winding of $Hx$, which is $+1$ for positive determinant and $-1$ for negative determinant: a positive-determinant matrix deforms to $I$, a negative one to $\operatorname{diag}(1,-1)$. Taylor's formula makes a definite Hessian a strict extremum and an indefinite Hessian a saddle, giving the asserted cases.
\end{proof}

\begin{theorem}\label{thm:morseexists}
Every compact smooth surface without boundary carries a Morse function. In particular it has a smooth tangent field with finitely many zeros; if it is oriented, its Euler characteristic above is defined.
\end{theorem}

\begin{proof}
The generic-height construction preceding the definition of Euler characteristic proves the statement: its exceptional parameter set is null, and the resulting nondegenerate critical points are isolated and form a closed subset of the compact surface, hence are finite. The gradient vanishes exactly there. That construction did not use orientation, so it applies also to nonorientable compact surfaces. On the empty surface the assertion is vacuous. For an oriented surface it supplies the field required by the definition of $\chi$.
\end{proof}

\begin{corollary}\label{cor:chimorse}
A Morse function on a compact oriented surface without boundary satisfies
$\chi(M)=\#\{\text{minima}\}-\#\{\text{saddles}\}+\#\{\text{maxima}\}$.
\end{corollary}

\begin{proof}
At each critical point the symmetric invertible $2\times2$ Hessian has either two positive eigenvalues, two negative ones or opposite signs. Taylor's formula gives respectively a strict minimum, a strict maximum or a saddle. The gradient indices are $+1,+1,-1$ by the lemma. Sum them over the finite critical set and use the definition of $\chi$.
\end{proof}

\begin{example}\label{ex:chiexamples}
The sphere's height function has a nondegenerate minimum and maximum, so $\chi(S^2)=2$ and its unit-radius curvature integral is $4\pi$. The torus has a nowhere-zero coordinate field $\partial_{\theta_1}$, giving $\chi(T^2)=0$. Therefore any metric on it has either identically zero curvature or curvature of both signs: a nonzero continuous function of only one sign could not integrate to zero.

For a standard genus-$g$ handle surface, one can choose a height with one minimum, one maximum and two saddles per handle; the count gives $2-2g$. The usual smooth classification of closed oriented surfaces by genus is quoted here for this interpretation, not proved by the height sketch. Chapter~6 will recover the same Euler characteristic from homology and the alternating count $V-E+F$ of a triangulation. Those later descriptions are comparisons, not premises for the curvature proof.
\end{example}

\begin{remark}
A nowhere-zero field on $S^2$ would make its index sum zero, contradicting $\chi=2$; this reproves the two-dimensional hairy-ball theorem. The general index-sum theorem is Poincar\'e--Hopf, with the higher-dimensional Euler characteristic defined by topology. We quote that extension: it holds for tangent fields with isolated zeros on compact manifolds without boundary. Chern's 1944 paper proves the even-dimensional Gauss--Bonnet formula by transgression on the unit tangent bundle. The surface proof above is the concrete case needed here.
\end{remark}

\section{The Laplacian of a surface}

The metric gives another operator on functions. On $1$-forms, the \emph{Hodge star} turns an oriented orthonormal coframe by a right angle: $\ast\theta^1=\theta^2$, $\ast\theta^2=-\theta^1$. This rule commutes with rotations, so it is frame independent. On $2$-forms set $\ast(h\,dA)=h$. Define the \emph{Laplace--Beltrami operator} by $\Delta f=\ast d\ast df$.

In Euclidean coordinates, $\ast df=f_xdy-f_ydx$ and $d\ast df=(f_{xx}+f_{yy})dx\wedge dy$. Thus this convention agrees with the usual Laplacian. Reversing orientation changes the two relevant star signs together and leaves $\Delta$ unchanged; it depends only on the metric and commutes with all isometries, including orientation-reversing ones.

\begin{proposition}[Green's formula]\label{prop:green}
On a compact oriented surface without boundary,
$\int_M g\Delta f\,dA=-\int_M\langle dg,df\rangle dA$ for smooth real $f,g$. Consequently $-\Delta$ is symmetric and nonnegative. Every eigenvalue $\lambda$ in $\Delta f=-\lambda f$, $f\ne0$, is real and nonnegative, and on a connected surface its zero eigenfunctions are exactly the constants.
\end{proposition}

\begin{proof}
Writing the two differentials in an orthonormal coframe gives $dg\wedge\ast df=\langle dg,df\rangle dA$. Hence
\[
d(g\ast df)=\langle dg,df\rangle dA+g\Delta f\,dA.
\]
Stokes makes the left integral zero, proving the formula. Interchanging $f,g$ proves symmetry, and taking $g=f$ proves nonnegativity. For a possibly complex eigenfunction, apply the bilinear identity to real and imaginary parts to obtain
$\lambda\int|f|^2dA=\int|df|^2dA$. Its nonzero positive denominator proves that $\lambda$ is real and nonnegative. If it is zero, the nonnegative continuous integrand vanishes, so $df=0$. A connected manifold is path connected, making $f$ constant; constants plainly satisfy $\Delta f=0$.
\end{proof}

\begin{theorem}[the spectrum of the square flat torus]\label{thm:torusspectrum}
On $T^2=\R^2/\Z^2$ with metric $dx^2+dy^2$, the eigenvalues of $-\Delta$ are $4\pi^2n$ for integers $n=k_1^2+k_2^2$, $k_1,k_2\in\Z$. Their multiplicities are
$r_2(n)=\#\{(k_1,k_2)\in\Z^2:k_1^2+k_2^2=n\}$. The functions $e_k(x)=e^{2\pi i\langle k,x\rangle}$ form an orthonormal eigenbasis of $L^2(T^2)$.
\end{theorem}

\begin{proof}
Differentiation gives $\Delta e_k=-4\pi^2|k|^2e_k$. Integrating the product of two modes over the unit square gives zero unless their two integer frequencies agree, and one if they do; this is the product of the one-variable orthogonality integrals.

Finite trigonometric sums form a self-conjugate algebra containing constants and separating points of the compact torus. Stone--Weierstrass, Volume~7, Theorem~4.19, gives uniform density in continuous functions; the latter are dense in $L^2$. Thus these orthonormal modes form a complete basis.

For a smooth eigenfunction, Green's identity with the conjugate mode gives
$(\lambda-4\pi^2|k|^2)\hat f(k)=0$ for every $k$. If no integer pair has the required norm, all coefficients vanish, contradicting $f\ne0$. Otherwise only the finitely many frequencies on that lattice circle remain. Completeness makes $f$ their linear combination; conversely every such combination is an eigenfunction. This proves precisely the stated multiplicities, including $r_2(0)=1$ for constants.
\end{proof}

For example, $n=1$ has four pairs $(\pm1,0),(0,\pm1)$, and $n=2$ has four pairs $(\pm1,\pm1)$. The counting function $N(\lambda)$, including multiplicity and the zero eigenvalue, is the number of lattice points in a disk of radius $R=\sqrt\lambda/(2\pi)$. Unit squares centred on those points put its area between the areas of disks with radii $\max(0,R-\sqrt2/2)$ and $R+\sqrt2/2$. Consequently $N(\lambda)=\lambda/(4\pi)+O(\sqrt\lambda+1)$ and in particular $N(\lambda)\sim\lambda/(4\pi)$ as $\lambda\to\infty$.

\begin{remark}
Weyl's law, quoted beyond this example, relates high eigenvalue counts to area for bounded plane domains under the usual regularity and boundary conditions. Weyl announced his result in 1911. Kac's 1966 question ``Can one hear the shape of a drum?'' asks how much geometry the spectrum determines. The error in the torus count connects to the Gauss circle problem of Volume~6, Chapter~2. We also quote the unit sphere spectrum $\ell(\ell+1)$ with multiplicity $2\ell+1$, obtained from spherical harmonics, and Milnor's 1964 example of nonisometric flat $16$-tori with equal spectra. None of these extensions is used to establish the square-torus calculation.
\end{remark}

\section{What we proved, and what comes next}

On the sphere, the curvature is $1$ and its integral is $4\pi$. On a torus, the integral must be zero even after the metric changes. These are constraints on every metric, rather than just calculations for familiar embeddings. The moving-frame definition makes the constraint intrinsic; the index sum explains why zeros of a vector field carry the same total.

For a second reading, follow the two uses of a derivative: $d\omega$ measures curvature, and the derivative of a Morse function supplies a field whose indices can be counted. The flat-torus modes then show a different way geometry meets arithmetic: a frequency pair's squared length is an eigenvalue divided by $4\pi^2$.

Euler characteristic is one number. The next chapter keeps the loops themselves, up to deformation, and gives them a group operation. This finer invariant distinguishes coverings such as $\R\to S^1$ and explains what changes when a loop cannot be contracted.

\section{Exercises}

\begin{figure}[ht]
\centering
\fig{ch04-girard}
\caption{A spherical triangle $ABC$ on the unit sphere, with the three great circles that carry its sides. Its angles add up to more than $\pi$; the excess is its area.}
\label{fig:girard}
\end{figure}

\begin{figure}[ht]
\centering\fig{v9-ch04-stereographic-length}
\caption{For the spherical metric $4(dx^2+dy^2)/(1+r^2)^2$, the coordinate circle of radius $r$ has length $4\pi r/(1+r^2)$. Its greatest length is $2\pi$ at $r=1$; very large coordinate circles approach the omitted pole and shrink in spherical length.}
\label{fig:v9-ch04-stereographic-length}
\end{figure}

\begin{figure}[ht]
\centering\fig{v9-ch04-frequency-circles}
\caption{Each lattice point is one square-torus Fourier mode. The circles $|k|^2=1,2,4$ each contain four lattice points, hence give eigenvalues $4\pi^2$, $8\pi^2$, $16\pi^2$ of multiplicity four. The origin gives the one-dimensional zero eigenspace.}
\label{fig:v9-ch04-frequency-circles}
\end{figure}

\begin{exercise}[$\star$]\label{ex:girard}
A spherical triangle here is a nondegenerate convex triangle on the unit sphere with minor great-circle sides, contained in an open hemisphere. A \emph{lune} of angle $\alpha$ is bounded by two half great circles joining antipodal points.
(a) Prove its area is $2\alpha$.
(b) Extend the triangle's sides to great circles. Show that the three containing lunes and their antipodal partners cover the triangle and its antipode three times each and every other point once, apart from boundaries.
(c) Deduce the triangle area $\alpha+\beta+\gamma-\pi$.
(d) Estimate the excess in arcseconds for area $2.9\cdot10^3$ km$^2$ on a sphere of radius $6371$ km.
\end{exercise}

\begin{exercise}[$\star$]\label{ex:revolution}
Rotate a unit-speed profile $(r(s),z(s))$, with $r>0$, around its vertical axis.
(a) Derive the metric $ds^2+r(s)^2d\vartheta^2$ and curvature $K=-r''/r$.
(b) For the torus made from a circle of radius $a>0$ whose centre is distance $R>a$ from the axis, compute $K$, its sign regions and its total curvature integral.
\end{exercise}

\begin{exercise}[$\star\star$]\label{ex:pseudosphere}
For $s>0$, set $r(s)=e^{-s}$.
(a) Find every smooth unit-speed profile $(r(s),z(s))$ of curvature $K=-1$ after revolution, including reflection and translation choices. Verify the tractrix property that its tangent segment to the axis has length $1$.
(b) Describe its local hyperbolic coordinates $x=\vartheta$, $y=e^s$ and the global identification $x\sim x+2\pi$. Explain why a real angle branch is needed for a map to the half-plane.
(c) Compute its area and compare with a unit hemisphere. For the upward profile, also find the volume of the horn bounded by the surface, its axis and its base disk.
\end{exercise}

\begin{exercise}[$\star$]\label{ex:chisphere}
(a) For every integer $n\ge1$, compute the indices at $0$ of the plane vector fields $z^n$ and $\bar z^n$.
(b) On $S^2\setminus\{N\}$ take the constant vector field $\partial_x$ in the complex stereographic chart $z$. In the chart $w=1/z$ about $N$, find its vector components and show that it extends smoothly with a single zero of index $2$. Deduce $\chi(S^2)=2$ once again.
\end{exercise}

\begin{exercise}[$\star$]\label{ex:positivepoint}
Let $S\subset\R^3$ be a nonempty compact smooth surface without boundary, and let $p$ maximize $|x|$ on it, with $R=|p|$. Prove $R>0$ and $K(p)\ge1/R^2$. Deduce that such a surface cannot have $K\le0$ everywhere, and that an embedded torus must have a point of negative curvature as well. Explain why the boundary hypothesis matters.
\end{exercise}

\begin{exercise}[$\star\star$]\label{ex:geodlength}
Let $U=\exp_p(D)$ for a tangent disk $D$ of radius $\rho$ on which $\exp_p$ is a diffeomorphism.
(a) Use the Gauss lemma, proved above, that the polar metric is $dr^2+G(r,\vartheta)d\vartheta^2$, to prove that every piecewise smooth curve from $p$ to $\exp_p(v)$, $|v|<\rho$, has length at least $|v|$, even if it leaves $U$. Describe equality up to monotone reparametrization, including $v=0$.
(b) Compute the exponential map explicitly at the sphere's North Pole and at $i$ in the half-plane. Derive their polar metrics locally, without using Chapter~9.
\end{exercise}

\begin{exercise}[$\star\star$]\label{ex:v9-ch04-conformal-chart}
For a smooth $u$ on a planar open set, equip it with metric $e^{2u}(dx^2+dy^2)$.
(a) Find its connection form in the coframe $e^u dx,e^u dy$ and prove $K=-e^{-2u}(u_{xx}+u_{yy})$.
(b) Check $K=-1$ for $u=-\log y$, $y>0$.
(c) Check $K=1$ for $u=\log2-\log(1+x^2+y^2)$, and compute the metric length of its coordinate circle of radius $r$. Explain why this length tends to zero as $r\to\infty$.
\end{exercise}

\begin{exercise}[$\star\star$]\label{ex:v9-ch04-conformal-change}
On a compact oriented Riemannian surface without boundary, replace $g$ by $\tilde g=e^{2u}g$, for a smooth real function $u$.
(a) Prove that the star on $1$-forms is unchanged, $d\tilde A=e^{2u}dA$, and $\tilde\omega=\omega+\ast du$ in correspondingly scaled frames.
(b) Deduce $\tilde K=e^{-2u}(K-\Delta u)$ and show directly that the total curvature is unchanged.
(c) If $g$ is flat on a torus, can a conformal change make its curvature strictly positive everywhere? If the new curvature is everywhere nonnegative, what must happen to $u$?
\end{exercise}

\begin{exercise}[$\star\star$]\label{ex:v9-ch04-rectangular-spectrum}
For $a,b>0$, consider $\R^2/(a\Z\times b\Z)$ with the Euclidean metric.
(a) Find a normalized Fourier eigenbasis and all multiplicities of its Laplacian.
(b) Find its first positive eigenvalue and multiplicity, distinguishing $a=b$ from $a\ne b$.
(c) Compare the unit square torus with the rectangle $a=2,b=1/2$. They have equal area; do they have equal spectra? Verify that both eigenvalue counting functions have leading term $ab\lambda/(4\pi)$ by a lattice-cell area estimate.
\end{exercise}

\begin{exercise}[$\star\star$]\label{ex:v9-ch04-split-index}
Let $f_\varepsilon(x,y)=x^3/3-xy^2-\varepsilon x$ and use the Euclidean metric.
(a) Compute its gradient and the index of its zero at $0$ when $\varepsilon=0$. Why is this not covered by the nondegenerate critical-point lemma?
(b) For either sign of $\varepsilon\ne0$, find all critical points and classify them by their Hessians.
(c) On a radius-$r$ circle with $|\varepsilon|<r^2$, show that the gradient has the same winding as at $\varepsilon=0$, and compare with the sum of the indices in the disk. This is a local example of a degenerate zero splitting without changing total index.
\end{exercise}

\section{Hints}

\begin{hint}{ex:girard}
Use spherical coordinates with the lune's endpoints as poles. For the coverage count, the three great circles give eight regions; each lune contains the triangle and one neighboring region.
\end{hint}
\begin{hint}{ex:revolution}
Compute the two coordinate tangent vectors and their dot products. Use the orthogonal-coordinate formula with $a=1,b=r(s)$.
\end{hint}
\begin{hint}{ex:pseudosphere}
Unit speed forces $z'^2=1-e^{-2s}>0$. Continuity fixes the sign of $z'$ on $s>0$. The angular variable lives modulo $2\pi$.
\end{hint}
\begin{hint}{ex:chisphere}
Follow arguments on $re^{it}$. For the second part transform a velocity by $dw/dt=-z^{-2}dz/dt$, rather than treating a vector field as a scalar function.
\end{hint}
\begin{hint}{ex:positivepoint}
Differentiate $|c(t)|^2$ twice along any surface curve through the farthest point. Relate its normal acceleration to the symmetric map $dN_p$.
\end{hint}
\begin{hint}{ex:geodlength}
Speed in polar coordinates dominates $|r'|$. A path leaving the chart crosses every smaller radius circle. In the half-plane, use the Cayley map only as an explicit coordinate change and calculate its metric directly.
\end{hint}
\begin{hint}{ex:v9-ch04-conformal-chart}
Put $a=b=e^u$ in Proposition~\ref{prop:orthocurv}. On a radius-$r$ circle, $u$ is constant in part (c).
\end{hint}
\begin{hint}{ex:v9-ch04-conformal-change}
Differentiate $e^u\theta^1,e^u\theta^2$ and compare the structure equations. Use $d\ast du=(\Delta u)dA$ and Green's formula with a constant test function.
\end{hint}
\begin{hint}{ex:v9-ch04-rectangular-spectrum}
Modes are $e^{2\pi i(kx/a+ly/b)}$. The shortest nonzero frequency lies along a longest side. For the count, surround each lattice point by its fixed rectangular frequency cell.
\end{hint}
\begin{hint}{ex:v9-ch04-split-index}
The complex gradient is $\bar z^2-\varepsilon$. On $|z|=r$ its unperturbed magnitude is $r^2$, giving a nonvanishing straight-line homotopy when $|\varepsilon|<r^2$.
\end{hint}

\section{Solutions}

\begin{solution}{ex:girard}
\idea Lune area is proportional to angle, and its overlap count isolates the triangle's excess.
\writeup (a) With the antipodal endpoints as poles, integrate the spherical area element $\sin\phi\,d\phi\,d\vartheta$ over $0\le\phi\le\pi$ and an angular interval of length $\alpha$. The result is $2\alpha$.
(b) The three distinct great circles cut the sphere into the triangle $T$, its antipode $T'$, and six remaining triangles. The lune at each vertex contains $T$ and the adjacent triangle across the opposite side; its antipodal partner contains $T'$ and the antipode of that adjacent triangle. The three choices account for all six other regions once. Thus $T,T'$ are counted three times and everything else once, apart from the great-circle boundaries of area zero.
(c) Add the six lune areas:
$4(\alpha+\beta+\gamma)=4\pi+4|T|$. This gives $|T|=\alpha+\beta+\gamma-\pi$.
(d) Scaling sphere lengths by $R$ scales area by $R^2$, so the excess is $2900/6371^2\approx7.145\cdot10^{-5}$ radians. Multiplication by $180\cdot3600/\pi$ gives approximately $14.74$ arcseconds, or about fifteen. The stated side lengths give the quoted area scale by Heron's formula; this numerical estimate illustrates the exact area--excess relation.
\end{solution}

\begin{solution}{ex:revolution}
\idea Orthogonal meridians and parallels make the curvature formula a second derivative of radius.
\writeup (a) The tangent vectors are $(r'\cos\vartheta,r'\sin\vartheta,z')$ and $(-r\sin\vartheta,r\cos\vartheta,0)$. Their squared lengths are $1,r^2$ and their dot product is zero. Thus $a=1,b=r$ and $K=-r''/r$.
(b) Use $r(s)=R+a\cos(s/a)$, $z(s)=a\sin(s/a)$ for one period $0\le s\le2\pi a$. Then
\[
K=\frac{\cos(s/a)}{a(R+a\cos(s/a))},\qquad dA=r(s)ds\,d\vartheta.
\]
Since $R>a$, its sign is the cosine's sign: positive on the outer half, negative on the inner half, zero on the top and bottom circles. Its integral is $2\pi\int_0^{2\pi a}\cos(s/a)ds/a=0$, as required for a torus.
\end{solution}

\begin{solution}{ex:pseudosphere}
\idea The prescribed radius fixes the metric, while unit speed leaves exactly a height translation and a reflection choice.
\writeup (a) Since $r''=r$, the curvature is $-1$. Unit speed requires $z'^2=1-e^{-2s}$, strictly positive for $s>0$. Smoothness on this connected interval makes its sign constant, so all profiles are
\[
z(s)=z_0\pm\int_0^s\sqrt{1-e^{-2u}}du
=z_0\pm\left(\operatorname{arcosh}(e^s)-\sqrt{1-e^{-2s}}\right).
\]
Their tangent vector is $(-r,z')$ of length $1$. Starting from $(r,z)$, advancing a parameter distance $1$ along that tangent line reaches the axis because its radial coordinate becomes $r-r=0$. The tangent segment therefore has length $1$, the tractrix property.
(b) The metric is $ds^2+e^{-2s}d\vartheta^2$. On any angle branch, $x=\vartheta,y=e^s$ gives $(dx^2+dy^2)/y^2=ds^2+e^{-2s}d\vartheta^2$. There is no globally real-valued angle on the circular parallels. The angular cover $s>0,\vartheta\in\R$ maps isometrically to $y>1$; the actual pseudosphere is its quotient by $x\mapsto x+2\pi$. Thus the local coordinate maps, not a single global map to the half-plane, express the isometry. The edge at $s=0$ also prevents this from realizing the complete hyperbolic plane; Hilbert's theorem rules out such a complete smooth realization in $\R^3$.
(c) Area is $\int_0^{2\pi}\int_0^\infty e^{-s}ds\,d\vartheta=2\pi$, the area of a unit hemisphere. For the upward profile the horn volume is
$\pi\int_0^\infty e^{-2s}\sqrt{1-e^{-2s}}ds
=\pi\int_0^1u\sqrt{1-u^2}du=\pi/3$,
half the unit hemisphere volume. Translation or reflection does not change either measurement.
\end{solution}

\begin{solution}{ex:chisphere}
\idea A field's vector components acquire a derivative factor when the coordinate changes.
\writeup (a) On $z=re^{it}$ the fields have arguments $nt$ and $-nt$, giving indices $n,-n$. The condition $n\ge1$ makes the origin an actual isolated zero; $n=0$ would be a nonzero constant field.
(b) The constant field has flow $z\mapsto z+t$. In the other chart its flow is $w\mapsto w/(1+tw)$, with initial velocity $-w^2$. It therefore extends smoothly at $w=0$ with exactly one zero, whose index is $2$; the fixed factor $-1$ adds only a constant angle. There are no zeros in the original constant-field chart, so the sphere's total index is $2$.
\end{solution}

\begin{solution}{ex:positivepoint}
\idea Maximal distance bounds every normal curvature by the tangent enclosing sphere.
\writeup If $R=0$ the entire surface would be a point, impossible for a nonempty smooth two-dimensional submanifold. At the maximum of $h(x)=|x|^2$, the derivative along $S$ vanishes, so $p$ is normal to its tangent plane. Choose $N(p)=p/R$. For a unit tangent $v$, take a smooth surface curve $c$ with $c(0)=p,c'(0)=v$. The interior maximum gives
$0\ge(|c|^2)''(0)=2+2\langle c''(0),p\rangle$, hence $\langle c'',N\rangle\le-1/R$.
Differentiating $\langle c',N(c)\rangle=0$ gives $\langle dN_p(v),v\rangle\ge1/R$. Since $dN_p$ is symmetric, both its eigenvalues are at least $1/R$, and $K(p)\ge1/R^2$.
Thus curvature cannot be nonpositive everywhere. On a torus its integral is zero, so positivity on a neighborhood of $p$ forces negativity elsewhere. A flat disk, with boundary, is a counterexample to the first conclusion if the boundary hypothesis is omitted: its farthest points need not be interior critical points of $h$.
\end{solution}

\begin{solution}{ex:geodlength}
\idea Radial displacement cannot exceed travelled length; trying to leave the chart costs more than the target radius.
\writeup (a) If $v=0$, the constant path has length zero and is the only path with that length. Suppose $0<|v|<\rho$. A path leaving $U$ must cross the boundary of every smaller compact normal disk of radius $\rho'<\rho$. Before that crossing its length is at least $\rho'$ by $|c'|\ge|r'|$. Let $\rho'\to\rho$; the total length is at least $\rho>|v|$.
If the path stays in $U$, consider the part after its last visit to $p$; that visit exists because the inverse image of $p$ is closed in the compact parameter interval. Away from that instant the polar metric gives $|c'|\ge|r'|$. Integrate from a slightly later time and pass to the limit to obtain length at least $|v|$. Equality forces any earlier part to be constant, the angular derivative to vanish on the remaining smooth pieces, and the radial derivative to be nonnegative. These are precisely monotone reparametrizations of the radial geodesic, including pauses.
(b) At the unit sphere's North Pole $p$, for a unit tangent $\xi$, the curve $\cos r\,p+\sin r\,\xi$ has unit speed and acceleration normal to the sphere, so its covariant acceleration is zero. Uniqueness gives $\exp_p(r\xi)=\cos r\,p+\sin r\,\xi$. Differentiating in $r$ and direction angle gives $dr^2+\sin^2r\,d\vartheta^2$ for $0<r<\pi$.

For the half-plane, set $w=(z-i)/(z+i)$, with inverse $z=i(1+w)/(1-w)$. Direct differentiation gives
\[
\frac{|dz|^2}{(\operatorname{Im}z)^2}=\frac{4|dw|^2}{(1-|w|^2)^2}.
\]
Put $w=\tanh(r/2)e^{i\vartheta}$. Substitution gives $dr^2+\sinh^2r\,d\vartheta^2$. Its coframe $dr,\sinh r\,d\vartheta$ has connection form $\cosh r\,d\vartheta$, so a constant-angle radial unit tangent has zero covariant derivative. These curves are therefore geodesics; the inverse Cayley expression extends smoothly at $r=0$, with initial unit vector $i e^{i\vartheta}$ at $i$. Consequently
\[
\exp_i(ri e^{i\vartheta})
=i\frac{1+\tanh(r/2)e^{i\vartheta}}{1-\tanh(r/2)e^{i\vartheta}}.
\]
This proves the requested local polar formula using only this chapter's connection and explicit coordinates. It does not assume the distance or isometry classification of Chapter~9.
\end{solution}

\begin{solution}{ex:v9-ch04-conformal-chart}
\idea Equal scale factors in two perpendicular directions turn the curvature formula into a Euclidean Laplacian.
\writeup (a) With $a=b=e^u$, we get $a_y/b=u_y$ and $b_x/a=u_x$. Thus $\omega=-u_y dx+u_xdy$ and $K=-e^{-2u}(u_{xx}+u_{yy})$.
(b) For $u=-\log y$, the sum of second derivatives is $1/y^2$, while $e^{-2u}=y^2$, giving $K=-1$.
(c) For $u=\log2-\log(1+r^2)$, direct differentiation gives $u_{xx}+u_{yy}=-4/(1+r^2)^2$ and $e^{-2u}=(1+r^2)^2/4$, so $K=1$. On the coordinate circle, $u$ is constant, so its length is its Euclidean circumference multiplied by $e^u$: $C(r)=4\pi r/(1+r^2)$. It is maximal at $r=1$, with value $2\pi$, and tends to zero at infinity. These are the lengths of latitude circles in stereographic sphere coordinates; approaching infinity approaches the missing pole, not an infinitely large circle on the sphere.
\end{solution}

\begin{solution}{ex:v9-ch04-conformal-change}
\idea Rescale the coframe and track exactly the extra turning term.
\writeup (a) The new orthonormal coframe is $e^u\theta^1,e^u\theta^2$. Applying star to either scaled member gives the scaled old star, proving that the star on $1$-forms is unchanged. Their wedge gives $d\tilde A=e^{2u}dA$. Write $du=u_1\theta^1+u_2\theta^2$, so $\ast du=-u_2\theta^1+u_1\theta^2$. Differentiating the scaled coframe and using the structure equations yields $\tilde\omega=\omega+\ast du$ by uniqueness.
(b) Differentiate:
\[
-\tilde K e^{2u}dA=d\tilde\omega
=d\omega+d\ast du=(-K+\Delta u)dA.
\]
Thus $\tilde K=e^{-2u}(K-\Delta u)$. Integrating and applying Green with constant test function gives $\int\Delta u\,dA=0$, hence $\int\tilde Kd\tilde A=\int KdA$ directly.
(c) For a flat starting torus, $\tilde K=-e^{-2u}\Delta u$. Strict positivity would force a strictly negative integral of $\Delta u$, impossible. Nonnegativity forces $\Delta u\le0$; its integral zero and continuity imply $\Delta u=0$. Green with test function $u$ gives $\int|du|^2=0$, so $u$ is constant on the connected torus. Only a constant scaling, still flat, is possible in that case.
\end{solution}

\begin{solution}{ex:v9-ch04-rectangular-spectrum}
\idea Period lengths divide frequencies, while the area sets the normalization.
\writeup (a) The normalized modes are $(ab)^{-1/2}e^{2\pi i(kx/a+ly/b)}$, $k,l\in\Z$. Their eigenvalues are $4\pi^2(k^2/a^2+l^2/b^2)$, and the multiplicity of a value is the number of integer pairs giving it. Orthogonality follows by integrating over $[0,a]\times[0,b]$; density follows from the same trigonometric-algebra argument as for the square, after rescaling coordinates. Zero has multiplicity one.
(b) Let $m=\max(a,b)$. Any nonzero integer pair gives $k^2/a^2+l^2/b^2\ge1/m^2$, attained by the two unit frequencies along a longest side. If $a\ne b$, these are the only minimizers, so multiplicity is two. If $a=b$, either axis gives a pair and multiplicity is four. The first positive eigenvalue is $4\pi^2/m^2$.
(c) The unit square has first positive eigenvalue $4\pi^2$ of multiplicity four. The area-one $2$-by-$1/2$ rectangle has first eigenvalue $\pi^2$ of multiplicity two, so their spectra differ. For the count, the frequency lattice has cell area $1/(ab)$ and centred rectangular cells of maximal radius $c=\tfrac12\sqrt{a^{-2}+b^{-2}}$. The union of cells centred at points of a radius-$R$ disk contains the disk of radius $\max(0,R-c)$ and is contained in the disk of radius $R+c$. Thus the number of points is $ab\pi R^2+O(R+1)$ for fixed $a,b$. Substituting $R=\sqrt\lambda/(2\pi)$ gives $N(\lambda)=ab\lambda/(4\pi)+O(\sqrt\lambda+1)$, the same leading term for both area-one examples.
\end{solution}

\begin{solution}{ex:v9-ch04-split-index}
\idea A boundary winding can stay fixed while a degenerate critical point separates into several ordinary ones.
\writeup (a) The gradient is $(x^2-y^2-\varepsilon,-2xy)$, or $\bar z^2-\varepsilon$. At $\varepsilon=0$ it has one zero at $0$ and argument $-2t$ on a counterclockwise circle, hence index $-2$. Its Hessian at $0$ vanishes, so the nondegenerate lemma does not apply.
(b) For $\varepsilon>0$, the zeros are $(\pm\sqrt\varepsilon,0)$; for $\varepsilon<0$, they are $(0,\pm\sqrt{-\varepsilon})$. There are no other zeros because $xy=0$. At any of them,
\[
H=\begin{pmatrix}2x&-2y\\-2y&-2x\end{pmatrix},\qquad
\det H=-4(x^2+y^2)=-4|\varepsilon|<0.
\]
Both points are nondegenerate saddles, each of index $-1$.
(c) On $|z|=r$, every field $\bar z^2-t\varepsilon$, $0\le t\le1$, has magnitude at least $r^2-|\varepsilon|>0$. This gives a homotopy of nonzero boundary vectors to the unperturbed field, so its winding remains $-2$. Both perturbed zeros lie inside since their radius is $\sqrt{|\varepsilon|}<r$, and their two indices add to $-2$. This verifies conservation in this example without assuming a general index-additivity theorem not yet proved here.
\end{solution}

\section*{Sources and further reading}
\addcontentsline{toc}{section}{Sources and further reading}

C.~F.~Gauss, ``Disquisitiones generales circa superficies curvas'', \emph{Comment.\ Soc.\ Reg.\ Sci.\ Gotting.\ Recent.} 6 (1828), presented on 8 October 1827;
English translation by A.~Hiltebeitel and J.~Morehead, \emph{General Investigations of Curved Surfaces} (Princeton, 1902; Dover reprint). On the survey and the triangle,
A.~I.~Miller, ``The myth of Gauss' experiment on the Euclidean nature of physical space'', \emph{Isis} 63 (1972), 345--348, and E.~Breitenberger, ``Gauss's
geodesy and the axiom of parallels'', \emph{Arch.\ Hist.\ Exact Sci.} 31 (1984), 273--289, which reads the evidence differently. W.~K.~B\"uhler,
\emph{Gauss: A Biographical Study} (Springer, 1981), Chapter~8, describes the Hanover survey and the heliotrope.

O.~Bonnet, ``M\'emoire sur la th\'eorie g\'en\'erale des surfaces'', \emph{J.\ \'Ec.\ Polytech.} 19 (1848), 1--146. S.-S.~Chern, ``A simple intrinsic proof of
the Gauss--Bonnet formula for closed Riemannian manifolds'', \emph{Ann.\ of Math.} 45 (1944), 747--752. H.~Hopf, ``Vektorfelder in $n$-dimensionalen
Mannigfaltigkeiten'', \emph{Math.\ Ann.} 96 (1927), 225--250. H.~Weyl, ``Das asymptotische Verteilungsgesetz der Eigenwerte linearer partieller
Differentialgleichungen'', \emph{Math.\ Ann.} 71 (1912), 441--479 (announced in 1911); M.~Kac, ``Can one hear the shape of a drum?'', \emph{Amer.\ Math.\ Monthly}
73 (1966), 1--23; J.~Milnor, ``Eigenvalues of the Laplace operator on certain manifolds'', \emph{Proc.\ Nat.\ Acad.\ Sci.\ USA} 51 (1964), 542.

The treatment by moving frames follows B.~O'Neill, \emph{Elementary Differential Geometry}, 2nd ed. (Academic Press, 2006), Chapters~6--7, and M.~do~Carmo,
\emph{Differential Forms and Applications} (Springer, 1994); for the classical approach, do~Carmo's \emph{Differential Geometry of Curves and Surfaces}
(Prentice-Hall, 1976). The existence of Morse functions follows V.~Guillemin and A.~Pollack, \emph{Differential Topology} (Prentice-Hall, 1974), Chapter~1, \S7.
